% Generated semantic source for AI Almanach v3.
% Edit v3 directly after initial generation; do not overwrite
% editorial changes by rerunning the converter casually.

% ==== Artificial Intelligence in FinTech =============================================================================
\chapter{Artificial Intelligence in FinTech}

\chapterlead{How can AI support financial decisions under asymmetric risk?}{Onboard}{Assess}{Transact}{Govern}{fintech}

% ============================================================
\label{ch:fintech}

\begin{learningobjectives}
After working through this chapter you should be able to:
\begin{itemize}
  \item compute precision and recall for a fraud detector at realistic base
        rates and explain why accuracy is meaningless there;
  \item define point-in-time correctness and name the leak that backtests
        create most often;
  \item explain why a model that predicts returns can still lose money after
        costs;
  \item identify a proxy variable that reintroduces a protected attribute into
        a credit decision;
  \item state what a regulator requires beyond predictive performance.
\end{itemize}
\end{learningobjectives}

\begin{plainlanguage}
Finance is where a small edge, applied often, becomes real money --- and where
a small methodological error, applied often, becomes a real loss. The
mathematics is mostly familiar from earlier chapters. What is different is that
the data moves, an adversary is optimising against you, the decisions are
legally consequential, and the thing you are predicting reacts to your having
predicted it.
\end{plainlanguage}

\begin{engineernote}
Two properties make financial machine learning unlike the tabular problems in
\cref{ch:machine-learning}. First, the data is non-stationary by construction:
a pattern that persisted for three years may vanish the month you deploy,
because someone else found it too. Second, every backtest is a measurement of a
past that cannot be re-run, so the only real test is forward and out of sample.
Treat any backtest as a hypothesis, never as evidence of profit.
\end{engineernote}

\section{Concepts of FinTech and Artificial Intelligence}
% ============================================================

% ---- 10.1.1 Introduction to FinTech ----------------------------------------------------
\subsection{Introduction to FinTech}


\begin{v3topic}{What is FinTech?}
    \textbf{FinTech} (Financial Technology) refers to the use of innovative technology to design, deliver, and improve financial products and services\textsuperscript{\cite[][p.~3]{kardellBuildFinancialSoftware2025}}.
    FinTech spans a broad spectrum: from mobile payment apps and digital lending platforms to algorithmic trading, blockchain, and AI-powered advisory services.
    The term gained widespread use after the 2008 financial crisis as startups challenged incumbents with leaner, data-driven business models.

\end{v3topic}
\begin{v3topic}{The FinTech Ecosystem}
    The FinTech ecosystem is a network of interdependent actors co-creating value around financial services.
\visualseparator
    \begin{tblr}{width=\linewidth,colspec={|Q[c,m,4cm]|X|}}
        \hline
        \textbf{Actor}        & \textbf{Role} \\
        \hline
        Incumbent banks       & Capital, trust, customer base, regulated infrastructure \\
        \hline
        FinTech startups      & Agile innovation, niche products, technology-first design \\
        \hline
        Regulators (BaFin, ECB, FCA) & Set rules, grant licenses, enforce compliance \\
        \hline
        BigTech (Apple, Google, Amazon) & Distribution power, data, embedded finance \\
        \hline
        Customers             & End users driving adoption and behavioral data \\
        \hline
    \end{tblr}
\captionof{table}{Actors in the FinTech ecosystem and what each contributes.}
\label{tab:fin-ecosystem}

\end{v3topic}
\begin{v3topic}{Open Banking and PSD2}
    The EU \textbf{Payment Services Directive 2} (PSD2, 2018) mandates that banks expose customer account data via standardized APIs to licensed third-party providers (TPPs)\textsuperscript{\cite[][p.~4]{kardellBuildFinancialSoftware2025}}.
\visualseparator
    \begin{itemize}
        \item \textbf{AISPs} (Account Information Service Providers): read-only access to account data
        \item \textbf{PISPs} (Payment Initiation Service Providers): initiate payments on behalf of customers
        \item Enables aggregation, personal finance management, and lending innovation
        \item Analogous open-banking frameworks adopted in the UK, Australia, and USA
    \end{itemize}

\end{v3topic}
\begin{v3topic}{Revolution in Financial Services}
    Three forces have disrupted traditional finance\textsuperscript{\cite[][pp.~1--2]{jansenMachineLearningAlgorithmic2020}}:
\visualseparator
    \begin{enumerate}
        \item \textbf{Digitization} of data — transaction records, market feeds, alternative data (social, satellite, IoT) are now machine-readable at scale
        \item \textbf{Democratization} of computing — cloud platforms, open-source ML libraries, and APIs lower barriers to entry
        \item \textbf{ML advances} — supervised, unsupervised, and reinforcement learning algorithms extract signals that were previously inaccessible
    \end{enumerate}

\end{v3topic}

% ---- 10.1.2 Applications in Banking and Finance ----------------------------------------
\subsection{Applications in Banking and Finance}


\begin{v3topic}{Neo-Banks and Digital-Only Banking}
    \textbf{Neo-banks} (Revolut, N26, Chime, Monzo) operate without physical branches, delivering the full current-account experience through a smartphone app\textsuperscript{\cite[][p.~5]{kardellBuildFinancialSoftware2025}}.
    They leverage cloud infrastructure and open APIs to offer lower fees, real-time notifications, and instant onboarding.
    Their data-native architecture feeds ML models for spend categorization, fraud scoring, and personalized financial coaching.

\end{v3topic}
\begin{v3topic}{Payments and the ACH Network}
    The \textbf{Automated Clearing House} (ACH) is the US electronic payment backbone, processing billions of credit and debit transactions annually\textsuperscript{\cite[][pp.~5--6]{kardellBuildFinancialSoftware2025}}.
    ACH files contain structured batch records (file header, batch header, entry detail, addenda, batch/file control) parsed by bank software.
    Alongside ACH, real-time payment rails (SEPA Instant, FedNow, SWIFT gpi) are reshaping global money movement toward 24/7 settlement.

\end{v3topic}
\begin{v3topic}{Wealth Management and Lending}
    AI transforms \textbf{wealth management} through automated portfolio construction, risk profiling, and goal-based planning at lower cost (robo-advisory, see Section 4).
    In \textbf{lending}, ML models score creditworthiness using alternative data (transaction history, social behavior, telco data) beyond traditional credit bureau scores, enabling faster decisions and broader financial inclusion.

\end{v3topic}
\begin{v3topic}{Financial Inclusion}
    An estimated 1.4 billion adults globally remain \emph{unbanked} (World Bank, 2021).
    FinTech addresses inclusion by offering mobile-money services (M-Pesa), micro-lending platforms, and credit scoring from non-traditional data.
    ML allows lenders to assess risk for thin-file customers who lack formal credit histories — reducing the dependence on collateral and expanding access.

\end{v3topic}

% ---- 10.1.3 Underlying Technologies ---------------------------------------------------
\subsection{FinTech Underlying Technologies}


\begin{v3topic}{Cloud Banking}
    Cloud platforms (AWS, Azure, GCP) supply banks with elastic compute, managed data pipelines, and ML services on demand\textsuperscript{\cite[][pp.~65--67]{kardellBuildFinancialSoftware2025}}.
    \textbf{Core banking-as-a-service} providers (Thought Machine, Mambu, Finxact) offer cloud-native ledger systems replacing decades-old COBOL mainframes.
    Key benefits: real-time data availability, rapid deployment of ML models, lower operational risk via managed SLAs.

\end{v3topic}
\begin{v3topic}{Blockchain, DLT and Smart Contracts}
    A \textbf{distributed ledger} (DLT) is a shared, append-only database replicated across nodes with cryptographic integrity.
    \textbf{Blockchain} chains blocks of transactions via hash pointers; Bitcoin (2009) proved the concept; Ethereum added programmability through \textbf{smart contracts} — self-executing code that enforces agreement terms without an intermediary.
    Financial use cases include: cross-border settlement (Ripple/XRP), tokenized securities, trade finance automation, and central bank digital currencies (CBDCs).

\end{v3topic}
\begin{v3topic}{ML and DL in Finance}
    Machine learning is the primary analytical engine of modern FinTech\textsuperscript{\cite[][pp.~1--3]{jansenMachineLearningAlgorithmic2020}}:
\visualseparator
    \begin{tblr}{width=\linewidth,colspec={|Q[c,m]|X|}}
        \hline
        \textbf{ML Type}       & \textbf{Finance Application} \\
        \hline
        Supervised             & Credit scoring, fraud classification, price forecasting \\
        \hline
        Unsupervised           & Customer segmentation, anomaly detection, eigenportfolios \\
        \hline
        Reinforcement Learning & Algorithmic trading agents, portfolio optimization \\
        \hline
        Deep Learning          & Time-series forecasting (LSTM), sentiment (BERT), OCR \\
        \hline
    \end{tblr}
\captionof{table}{Machine-learning families mapped to financial applications.}
\label{tab:fin-ml-applications}

\end{v3topic}
\begin{v3topic}{NLP for Financial Text}
    Financial institutions generate and consume vast quantities of text: earnings call transcripts, regulatory filings (SEC EDGAR, annual reports), news feeds, and social media.
    NLP methods extract structured signals from unstructured text\textsuperscript{\cite[][pp.~440--441]{jansenMachineLearningAlgorithmic2020}}: sentiment scores from analyst reports, topic models from SEC filings, named-entity recognition for event detection, and large language model (LLM) summaries of earnings calls.

\end{v3topic}

% ============================================================
\section{AI Applications in Financial Services}
% ============================================================

% ---- 10.2.1 Credit Scoring and Lending ------------------------------------------------
\subsection{Credit Scoring and Lending}


\begin{v3topic}{Traditional Credit Scorecards}
    Classical credit scoring assigns a numerical score to an applicant based on a weighted sum of features: payment history, amounts owed, length of credit history, new credit, and credit mix (FICO model)\textsuperscript{\cite[][p.~149]{jansenMachineLearningAlgorithmic2020}}.
    Scorecards were historically built with \textbf{logistic regression} on bureau data, offering interpretability and regulatory defensibility but limited predictive power for thin-file customers.

\end{v3topic}
\begin{v3topic}{ML-Based Credit Models}
    Modern lenders deploy gradient boosting (XGBoost, LightGBM) and deep learning models on enriched feature sets\textsuperscript{\cite[][pp.~365--366]{jansenMachineLearningAlgorithmic2020}}.
    Alternative data signals:
\visualseparator
    \begin{itemize}
        \item Transaction patterns (cash-flow regularity, merchant categories)
        \item Mobile metadata (device usage, call records — with consent)
        \item Psychographic questionnaire scores
        \item Social network behavior (in some jurisdictions)
    \end{itemize}
    Tree ensembles handle non-linear interactions and missing data naturally; SHAP values provide local explanations.

\end{v3topic}
\begin{v3topic}{Fairness and Disparate Impact}
    Algorithmic credit decisions can inadvertently encode historical discrimination if protected attributes (race, gender, age) correlate with training features\textsuperscript{\cite[][p.~716]{jansenMachineLearningAlgorithmic2020}}.
    \textbf{Disparate impact} (US Equal Credit Opportunity Act) occurs when a neutral criterion disproportionately disadvantages a protected class.
    Mitigation: \emph{pre-processing} (re-weighting training data), \emph{in-processing} (fairness-regularized loss), \emph{post-processing} (threshold equalization across groups), and individual explainability requirements under GDPR Article 22.

\end{v3topic}

% ---- 10.2.2 Algorithmic Trading -------------------------------------------------------
\subsection{Algorithmic Trading}


\begin{v3topic}{What is Algorithmic Trading?}
    \textbf{Algorithmic trading} uses computer programs to execute trading decisions automatically according to pre-specified rules or learned strategies\textsuperscript{\cite[][pp.~1--2]{jansenMachineLearningAlgorithmic2020}}.
    The goal of active investment management is to generate \textbf{alpha} — portfolio returns in excess of a benchmark.
    The \emph{fundamental law of active management} decomposes the information ratio (IR) as:
    \begin{equation}
        \mathrm{IR} \approx \mathrm{IC} \cdot \sqrt{N}
    \end{equation}
    where IC is the information coefficient (forecast skill) and $N$ is the breadth (number of independent bets).

\end{v3topic}
\begin{v3topic}{Alpha Factors and Feature Engineering}
    An \textbf{alpha factor} is any signal, attribute, or transformation of data that historically correlates with future asset returns\textsuperscript{\cite[][pp.~82--84]{jansenMachineLearningAlgorithmic2020}}.
    Classic factor families: \emph{momentum} (12-1 month return), \emph{value} (price-to-book ratio), \emph{quality} (return on equity), \emph{volatility} (low-vol anomaly), \emph{size} (small-cap premium).
    ML feature engineering extends factors to alternative data: NLP sentiment scores, satellite imagery activity, web traffic, credit-card transaction aggregates.

\end{v3topic}
\begin{v3topic}{Strategy Types}
\visualseparator
    \begin{tblr}{width=\linewidth,colspec={|Q[c,m]|X|}}
        \hline
        \textbf{Strategy}       & \textbf{Description} \\
        \hline
        Statistical arbitrage   & Pairs trading / cointegration; mean-reversion of spread \\
        \hline
        Market making           & Post limit orders on both sides; profit from bid-ask spread \\
        \hline
        Trend following (CTA)   & Long/short based on momentum signals across asset classes \\
        \hline
        High-frequency (HFT)    & Microsecond latency; exploit microstructure inefficiencies \\
        \hline
        Long-short equity       & ML signals rank stocks; long winners, short losers \\
        \hline
    \end{tblr}
\captionof{table}{Algorithmic trading strategies by mechanism. Each depends on a market condition that other participants are also trying to exploit.}
\label{tab:fin-strategies}

\end{v3topic}
\begin{v3topic}{RL for Trading}
    Deep reinforcement learning casts trading as a Markov Decision Process: the agent observes market state $s_t$ (price features, portfolio state), takes action $a_t$ (buy/hold/sell), and receives reward $r_t$ (risk-adjusted P\&L)\textsuperscript{\cite[][pp.~679--680]{jansenMachineLearningAlgorithmic2020}}.
    The agent learns a policy $\pi(a|s)$ that maximizes expected cumulative return.
    Key challenge: non-stationarity of financial markets makes generalization from historical training data difficult.

\end{v3topic}
\begin{v3topic}{Backtesting and Overfitting Risk}
    Backtesting simulates a strategy on historical data to estimate future performance\textsuperscript{\cite[][pp.~222--226]{jansenMachineLearningAlgorithmic2020}}.
    Critical pitfalls:
\visualseparator
    \begin{tblr}{width=\linewidth,colspec={|X|X|X|}}
        \hline
        \textbf{Pitfall}         & \textbf{Cause}                                    & \textbf{Remedy} \\
        \hline
        Look-ahead bias          & Using future data in feature construction          & Strict point-in-time data discipline \\
        \hline
        Survivorship bias        & Only including assets that survived to today        & Use delisted/bankrupt-company data \\
        \hline
        Overfitting (data mining) & Too many strategy parameters fit to noise          & Walk-forward validation, Deflated Sharpe \\
        \hline
        Transaction costs        & Ignoring slippage, commissions, market impact       & Model realistic costs per strategy \\
        \hline
    \end{tblr}
\captionof{table}{Backtest pitfalls with cause and remedy. Every one of them makes the backtest look better and the deployment worse.}
\label{tab:fin-backtest-pitfalls}

\end{v3topic}

% ---- 10.2.3 Chatbots and Customer Service --------------------------------------------
\subsection{Chatbots and Customer Service}


\begin{v3topic}{NLP-Powered Banking Assistants}
    Virtual banking assistants (Bank of America's Erica, Cleo, KAI by Kasisto) handle natural language queries about balances, transactions, and product recommendations\textsuperscript{\cite[][p.~14]{kardellBuildFinancialSoftware2025}}.
    They combine \textbf{intent recognition} (classify user goal: balance inquiry, payment initiation, fraud report), \textbf{slot filling} (extract entities: amount, date, account), and \textbf{dialogue management} (maintain context across turns).
    Large language models (GPT-4, Claude) are now being integrated for open-ended advisory conversations.

\end{v3topic}
\begin{v3topic}{Generative AI in Finance}
    Generative AI enables FinTech applications beyond rule-based chatbots\textsuperscript{\cite[][pp.~13--15]{kardellBuildFinancialSoftware2025}}:
\visualseparator
    \begin{itemize}
        \item \textbf{Code generation}: AI pair-programmers assist developers building ACH parsers, API handlers, and database schemas
        \item \textbf{Document summarization}: Automatic digests of lengthy regulatory filings and contracts
        \item \textbf{Report generation}: Narrative commentary on portfolio performance
        \item \textbf{Synthetic data}: Training data for fraud detection where real data is scarce or sensitive
    \end{itemize}

\end{v3topic}

% ---- 10.2.4 KYC and AML --------------------------------------------------------------
\subsection{KYC and AML}


\begin{v3topic}{Know Your Customer (KYC)}
    \textbf{KYC} is the regulatory obligation for financial institutions to verify the identity of customers before and during a business relationship.
    Steps: identity document verification (OCR + liveness detection), sanctions list screening (OFAC, UN), politically exposed persons (PEP) checks, and ongoing monitoring\textsuperscript{\cite[][pp.~378--379]{kardellBuildFinancialSoftware2025}}.
    ML accelerates KYC via facial recognition, document forgery detection (CNN classifiers), and fuzzy-matching algorithms for name screening.

\end{v3topic}
\begin{v3topic}{Anti-Money Laundering (AML)}
    \textbf{AML} programs detect and report suspicious transactions that may represent layering or integration of illicit funds.
    Traditional rule-based systems generate high false-positive rates ($>$95\%).
    ML approaches:
\visualseparator
    \begin{itemize}
        \item \textbf{Transaction monitoring}: anomaly detection flags structuring (smurfing), round-trip transfers, and mule accounts
        \item \textbf{Network analysis}: graph ML (GNN) uncovers money laundering rings via transaction flow topology
        \item \textbf{Risk scoring}: ensemble models produce customer risk scores updated in real time
        \item \textbf{RegTech}: automated SAR (Suspicious Activity Report) narrative generation
    \end{itemize}

\end{v3topic}

% ============================================================
\section{Fraud Detection in FinTech}
% ============================================================

% ---- 10.3.1 Types of Financial Fraud -------------------------------------------------
\subsection{Types of Financial Fraud}


\begin{v3topic}{Payment and Card Fraud}
    \textbf{Payment fraud} includes card-not-present (CNP) fraud (online transactions using stolen card data), counterfeit cards (skimming), account takeover (ATO), and authorized push payment (APP) scams where customers are socially engineered into authorizing transfers.
    Global card fraud losses exceeded \$33 billion in 2022 (Nilson Report).
    Real-time fraud scoring at the point of authorization (sub-100 ms) is the primary defense.

\end{v3topic}
\begin{v3topic}{Insurance and Identity Fraud}
    \textbf{Insurance fraud} ranges from exaggerated claims (soft fraud) to staged accidents (hard fraud), costing the US industry an estimated \$80 billion annually.
    \textbf{Identity fraud} involves using stolen or synthetic personal information (name, SSN, date of birth) to open accounts, apply for credit, or conduct transactions.
    \textbf{Synthetic identity fraud} — combining real and fabricated data — is the fastest-growing fraud type and is particularly difficult to detect with traditional checks.

\end{v3topic}
\begin{v3topic}{The Wirecard Case}
    \textbf{Wirecard AG} (Germany) was a DAX-listed FinTech payment processor that collapsed in June 2020 after disclosing that \texteuro{}1.9 billion in cash held in Philippine escrow accounts did not exist.
    Key fraud mechanisms: fictitious revenue from third-party acquiring (TPA) partners, inflated balance sheet, and auditor failures (EY signed off for years).
    Lessons for AI in finance: anomaly detection on financial statements (Benford's law, ratio analysis), graph-based counterparty risk, and the need for explainable AI to support human investigator judgment rather than replace it.

\end{v3topic}

% ---- 10.3.2 ML for Fraud Detection ---------------------------------------------------
\subsection{ML for Fraud Detection}


\begin{v3topic}{Anomaly Detection Approaches}
    \textbf{Unsupervised anomaly detection} identifies unusual patterns without labeled fraud examples\textsuperscript{\cite[][p.~150]{jansenMachineLearningAlgorithmic2020}}:
\visualseparator
    \begin{itemize}
        \item \textbf{Isolation Forest}: randomly partitions features; anomalies require fewer splits to isolate
        \item \textbf{Autoencoder}: reconstruction error flags transactions the network cannot compress faithfully
        \item \textbf{One-class SVM}: learns a decision boundary around normal behavior
        \item \textbf{DBSCAN}: density-based clustering; low-density points flagged as anomalies
    \end{itemize}

\end{v3topic}
\begin{v3topic}{Supervised Classification for Fraud}
    When labeled fraud cases are available, supervised classifiers achieve higher precision:
    Gradient boosting (XGBoost, LightGBM), random forests, and neural networks are state of the art\textsuperscript{\cite[][pp.~345--346]{jansenMachineLearningAlgorithmic2020}}.
    Key features: transaction amount, merchant category code (MCC), device fingerprint, geolocation distance from home, time-of-day, velocity (number of transactions per hour), and behavioral biometrics.

\end{v3topic}
\begin{v3topic}{Class Imbalance Problem}
    Fraud datasets are severely imbalanced — typically 0.1--1\% positive class.
    Naively training on raw data produces a model that predicts ``not fraud'' for everything with high accuracy but zero utility.
\visualseparator
    Remedies:
    \begin{itemize}
        \item \textbf{SMOTE} (Synthetic Minority Oversampling Technique): interpolates synthetic minority samples in feature space
        \item \textbf{Cost-sensitive learning}: assign higher misclassification cost to false negatives
        \item \textbf{Threshold tuning}: optimize the decision threshold on precision-recall curve, not accuracy
        \item \textbf{Ensemble methods}: combine multiple weak detectors to improve recall
    \end{itemize}

\end{v3topic}
\begin{v3topic}{Real-Time Fraud Scoring}
    Production fraud systems must return a decision within $<$100 ms at transaction authorization time.
    Architecture: a lightweight feature extraction layer reads recent transaction history from an in-memory store (Redis), passes features to a pre-trained ML model served via REST API, and returns a risk score.
    Scores above a threshold trigger step-up authentication or decline; scores in a gray zone route to human review queues.
    Model drift monitoring and frequent retraining cycles are essential as fraud patterns evolve rapidly.
    \factvisual{
        \textbf{A score is not yet a decision.}
        \begin{itemize}
            \item Low risk: approve with ordinary monitoring.
            \item Uncertain: request step-up authentication.
            \item High risk: decline or hold according to policy.
            \item Afterwards: capture outcomes and appeals for monitoring.
        \end{itemize}
    }{
        \begin{tikzpicture}[node distance=4mm]
            \node[almanach card,text width=24mm] (tx) {transaction};
            \node[almanach card,text width=24mm,below=of tx] (features)
                {live + historical features};
            \node[almanach accent card,text width=24mm,below=of features] (score)
                {risk score};
            \node[almanach card,text width=18mm,below left=7mm and 8mm of score] (approve)
                {approve};
            \node[almanach accent card,text width=18mm,below=7mm of score] (verify)
                {verify};
            \node[almanach card,text width=18mm,below right=7mm and 8mm of score] (hold)
                {hold / decline};
            \draw[almanach arrow] (tx) -- (features);
            \draw[almanach arrow] (features) -- (score);
            \draw[almanach arrow] (score) -- node[left,font=\tiny] {low} (approve);
            \draw[almanach arrow] (score) -- node[right,font=\tiny] {grey} (verify);
            \draw[almanach arrow] (score) -- node[right,font=\tiny] {high} (hold);
        \end{tikzpicture}
        \visualcaption{A production fraud model feeds a threshold policy with
          three operational paths, each requiring outcome monitoring.}
          {fig:fraud-scoring-policy}
    }

\end{v3topic}

% ============================================================
\section{Robo-Advisory}
% ============================================================

% ---- 10.4.1 Definition and Types -------------------------------------------------------
\subsection{Definition and Types of Robo-Advisors}


\begin{v3topic}{What is a Robo-Advisor?}
    A \textbf{robo-advisor} is a digital platform that provides automated, algorithm-driven financial planning and investment management services with minimal human intervention\textsuperscript{\cite[][p.~125]{jansenMachineLearningAlgorithmic2020}}.
    First-generation robo-advisors (Betterment 2010, Wealthfront 2011) offered low-cost passive ETF portfolios.
    Current platforms provide goal-based planning, tax optimization, and personalized risk profiling at management fees of 0.15--0.50\% vs.\ 1--2\% for traditional advisors.

\end{v3topic}
\begin{v3topic}{Pure vs.\ Hybrid Robo-Advisors}
\visualseparator
    \begin{tblr}{width=\linewidth,colspec={|Q[c,m]|X|X|}}
        \hline
                             & \textbf{Pure Robo} & \textbf{Hybrid / Bionic} \\
        \hline
        Advice channel        & Fully automated algorithm & Algorithm + human advisor access \\
        \hline
        Cost                  & Very low (0.15--0.30\%)   & Moderate (0.50--0.85\%) \\
        \hline
        Complexity handled    & Standard risk profiles    & Complex tax, estate, business needs \\
        \hline
        Examples              & Betterment, Wealthfront   & Vanguard Personal Advisor, Schwab \\
        \hline
    \end{tblr}
\captionof{table}{Pure against hybrid robo-advice. The human element is retained mainly for the cases where the constraint is behavioural rather than computational.}
\label{tab:fin-robo-types}

\end{v3topic}

% ---- 10.4.2 Portfolio Management -------------------------------------------------------
\subsection{Portfolio Management}


\begin{v3topic}{Modern Portfolio Theory (Markowitz)}
    Harry Markowitz (1952) showed that a rational investor maximizes expected return for a given level of risk (variance)\textsuperscript{\cite[][p.~127]{jansenMachineLearningAlgorithmic2020}}.
    For a portfolio of $N$ assets with weight vector $\mathbf{w}$, expected return $\boldsymbol{\mu}$, and covariance $\boldsymbol{\Sigma}$:
    \begin{equation}
        \min_{\mathbf{w}}\; \mathbf{w}^\top \boldsymbol{\Sigma}\, \mathbf{w} \quad \text{s.t.} \quad \mathbf{w}^\top \boldsymbol{\mu} = \mu^*, \; \mathbf{1}^\top \mathbf{w} = 1
    \end{equation}
\where{$\mathbf{w}$ is the vector of portfolio weights, $\boldsymbol{\Sigma}$ the
covariance matrix of asset returns, $\boldsymbol{\mu}$ their expected returns,
and $\mu^*$ the required return. The last constraint forces the weights to sum
to one --- the whole portfolio is invested.}
    The set of optimal portfolios traces the \emph{efficient frontier} in mean-variance space.

\end{v3topic}
\begin{v3topic}{Sharpe Ratio and Risk-Adjusted Performance}
    The \textbf{Sharpe ratio} measures risk-adjusted return\textsuperscript{\cite[][p.~122]{jansenMachineLearningAlgorithmic2020}}:
    \begin{equation}
        S = \frac{\mathbb{E}[r_p - r_f]}{\sigma_p}
    \end{equation}
    where $r_p$ is portfolio return, $r_f$ is the risk-free rate, and $\sigma_p$ is portfolio return volatility.
    Robo-advisors target a client's required Sharpe ratio by tuning the equity/bond allocation.
    Modern alternatives include Sortino ratio (downside deviation only) and Calmar ratio (return/max drawdown).

\end{v3topic}
\begin{v3topic}{Black-Litterman Model}
    The \textbf{Black-Litterman model} (1990) addresses the instability of Markowitz optimization under estimation error in $\boldsymbol{\mu}$\textsuperscript{\cite[][pp.~131--132]{jansenMachineLearningAlgorithmic2020}}.
    It blends market-implied equilibrium returns (derived from the market-cap-weighted CAPM) with investor views via a Bayesian update:
    \begin{equation}
        \boldsymbol{\mu}^* = [(\tau \boldsymbol{\Sigma})^{-1} + \mathbf{P}^\top \boldsymbol{\Omega}^{-1} \mathbf{P}]^{-1}[\,(\tau \boldsymbol{\Sigma})^{-1} \boldsymbol{\Pi} + \mathbf{P}^\top \boldsymbol{\Omega}^{-1} \mathbf{q}\,]
    \end{equation}
    where $\boldsymbol{\Pi}$ is the equilibrium return, $\mathbf{P}$ is the view matrix, $\mathbf{q}$ the view vector, and $\boldsymbol{\Omega}$ view uncertainty.

\end{v3topic}
\begin{v3topic}{Rebalancing and Tax-Loss Harvesting}
    \textbf{Automatic rebalancing} periodically restores portfolio weights to target allocations when drift (due to differential asset returns) exceeds a threshold.
    Trigger-based rebalancing (e.g., when any asset deviates $>$5\%) is more tax-efficient than calendar-based.
    \textbf{Tax-loss harvesting} (TLH) systematically sells positions at a loss to realize capital losses that offset taxable gains, then repurchases a correlated substitute to maintain market exposure — subject to wash-sale rules (30-day window in the US).
    Wealthfront claims 1.55--2.99\% after-tax return improvement via daily TLH.

\end{v3topic}

% ---- 10.4.3 Goal-Based Investing -------------------------------------------------------
\subsection{Goal-Based Investing}


\begin{v3topic}{Risk Profiling and Investor Preferences}
    Robo-advisors elicit \textbf{risk tolerance} through questionnaires covering time horizon, income stability, financial knowledge, and emotional reaction to losses (based on behavioral finance research).
    Risk scores map to model portfolios along a conservative-to-aggressive spectrum (e.g., 20/80 to 100/0 equity/bond allocation).
    Limitations: self-reported risk tolerance is volatile and may diverge from revealed preference during market crashes.

\end{v3topic}
\begin{v3topic}{Goal-Based Portfolio Modeling}
    Rather than optimizing a single portfolio, \textbf{goal-based investing} allocates separate sub-portfolios (``buckets'') to distinct financial goals: emergency fund (short horizon, capital preservation), home purchase (medium horizon), and retirement (long horizon, growth)\textsuperscript{\cite[][p.~125]{jansenMachineLearningAlgorithmic2020}}.
    Each sub-portfolio has its own risk budget, return objective, and time horizon.
    Monte Carlo simulation projects goal achievement probabilities under thousands of market scenarios, enabling advice such as ``increase monthly contributions by \$150 to reach 90\% confidence.''

\end{v3topic}

% ============================================================
\section{Trust, Ethics, and Regulation}
% ============================================================

% ---- 10.5.1 Bias and Fairness ----------------------------------------------------------
\subsection{Bias and Fairness in AI Finance}


\begin{v3topic}{Algorithmic Discrimination in Lending}
    ML credit models may perpetuate or amplify historical discrimination even without explicitly encoding protected attributes\textsuperscript{\cite[][p.~716]{jansenMachineLearningAlgorithmic2020}}.
    Proxy discrimination arises when correlated features (zip code, education level) carry demographic signal.
    \textbf{Disparate impact}: under the US Equal Credit Opportunity Act (ECOA), the 80\% rule states that if the acceptance rate for a protected group is less than 80\% of the highest group rate, the policy is suspect.

\end{v3topic}
\begin{v3topic}{Explainability Requirements}
    Regulations demand explainability for automated financial decisions:
\visualseparator
    \begin{itemize}
        \item \textbf{GDPR Art.~22}: right not to be subject to solely automated decisions; right to explanation
        \item \textbf{EU AI Act (2024)}: credit scoring and insurance pricing classified as high-risk AI — mandatory transparency, logging, human oversight, and bias monitoring
        \item \textbf{US FCRA}: adverse action notices must cite specific reasons for credit denial
        \item SHAP and LIME provide model-agnostic local explanations compatible with regulatory requirements
    \end{itemize}

\end{v3topic}

% ---- 10.5.2 Financial Regulation ------------------------------------------------------
\subsection{Regulation of AI in Financial Services}


\begin{v3topic}{MiFID II and Algorithmic Trading}
    The EU \textbf{Markets in Financial Instruments Directive II} (MiFID II, 2018) introduced specific requirements for algorithmic trading firms:
\visualseparator
    \begin{itemize}
        \item Algorithm registration and stress-testing obligations
        \item Kill-switch capability to halt trading immediately
        \item Annual self-assessment of trading algorithms
        \item Market-making algorithms must maintain quotes during normal conditions
        \item Tick data reporting and best-execution obligations
    \end{itemize}

\end{v3topic}
\begin{v3topic}{Regulation of Robo-Advisors}
    Robo-advisors providing investment advice are subject to financial regulations across jurisdictions:
\visualseparator
    \begin{tblr}{width=\linewidth,colspec={|Q[c,m]|X|}}
        \hline
        \textbf{Jurisdiction} & \textbf{Regime} \\
        \hline
        EU                    & MiFID II suitability assessment; ESMA Guidelines on Robo-Advice (2018) \\
        \hline
        Germany               & BaFin authorization as investment advisor; GDPR compliance \\
        \hline
        USA                   & SEC Registered Investment Advisor (RIA) status; fiduciary duty \\
        \hline
        UK                    & FCA authorization; Consumer Duty (2023) \\
        \hline
    \end{tblr}
\captionof{table}{Robo-advice regimes by jurisdiction. Verify against current primary sources before relying on any row.}
\label{tab:fin-robo-regulation}

\end{v3topic}

% ---- 10.5.3 Future Outlook -------------------------------------------------------------
\subsection{Future Outlook}


\begin{v3topic}{CBDCs and Programmable Money}
    \textbf{Central bank digital currencies} (CBDCs) are sovereign digital money issued directly by central banks.
    Over 130 countries are exploring CBDCs (BIS survey 2023).
    Programmable CBDCs can embed spending conditions (welfare payments restricted to food), automate tax withholding, and enable real-time cross-border settlement without correspondent banking intermediaries.
    AI systems will be critical for CBDC fraud monitoring and AML compliance at scale.

\end{v3topic}
\begin{v3topic}{DeFi and Embedded Finance}
    \textbf{Decentralized Finance} (DeFi) replicates financial services (lending, trading, derivatives) through smart contracts on public blockchains, without centralized intermediaries.
    Total value locked (TVL) in DeFi peaked at $\sim$\$180 billion in 2021.
    \textbf{Embedded finance} integrates financial services into non-financial platforms: Shopify Capital (merchant lending), Uber Money, and BNPL (buy now, pay later) at checkout.
    ML models are the intelligence layer enabling real-time underwriting at point of sale.

\end{v3topic}
\begin{v3topic}{ESG and Sustainable Finance}
    \textbf{ESG} (Environmental, Social, Governance) investing integrates non-financial risk factors into asset allocation.
    AI applications in sustainable finance\textsuperscript{\cite[][p.~715]{jansenMachineLearningAlgorithmic2020}}:
\visualseparator
    \begin{itemize}
        \item Satellite imagery and NLP to monitor corporate environmental compliance
        \item LLMs to extract ESG disclosures from annual reports and sustainability reports
        \item ML to construct ESG scores from unstructured data beyond agency ratings
        \item Climate risk stress-testing of loan portfolios (NGFS scenarios)
    \end{itemize}

\end{v3topic}

% ============================================================
%
\section{Deep Dive: Point-in-Time Financial Evaluation}
\begin{v3topic}{What Was Knowable Then?}
Reconstruct features, universe membership, prices, and filings as they were
available at the decision timestamp.
Later corrections, surviving firms, or revised labels create look-ahead and
survivorship bias even when rows are split by date.

\end{v3topic}
\begin{v3topic}{Backtest the Executable Policy}
Include transaction costs, spread, latency, market impact, turnover,
constraints, rejected orders, and capital limits.
Compare against a simple implementable strategy and test performance across
regimes rather than selecting a convenient interval.

\end{v3topic}
\begin{v3topic}{Decisions Need Reasons and Appeals}
For credit, fraud, or AML workflows, connect model output to threshold,
capacity, adverse-action explanation, human review, and appeal.
Measure calibration and error costs by relevant groups, plus queue burden and
time to resolution.
\end{v3topic}


\begin{workedexample}[Why a fraud detector with 99.9\,\% accuracy is unusable]
One million transactions, of which $0.1\,\%$ --- one thousand --- are
fraudulent. A model catches $800$ of them and raises $9{,}000$ false alarms.

\[
\begin{aligned}
  \text{Accuracy} &= \frac{800 + (999{,}000-9{,}000)}{1{,}000{,}000}
    = \frac{990{,}800}{1{,}000{,}000} = 99.08\,\%,\\[3pt]
  \text{Precision} &= \frac{800}{800+9{,}000} = \frac{800}{9{,}800}
    = 8.2\,\%,\\[3pt]
  \text{Recall} &= \frac{800}{1{,}000} = 80\,\% .
\end{aligned}
\]

\emph{The do-nothing baseline.} Approving everything scores
$99.9\,\%$ accuracy --- \emph{better} than the model.

\emph{The operational reading.} Of every twelve alerts a reviewer opens, eleven
are legitimate customers. At $9{,}800$ alerts and, say, four minutes each, that
is roughly $650$ analyst-hours per million transactions. The binding constraint
is not the model's quality but the review capacity, and the threshold must be
set from that capacity, not from an $F_1$ optimum.

\emph{The decision that follows.} Compare the expected cost of a missed fraud
against the cost of a false alarm --- including the customer whose legitimate
card was declined at a checkout. Those two numbers, not the ROC curve, choose
the threshold.
\end{workedexample}

\begin{cautionbox}[Point-in-time correctness, and the leak that flatters every backtest]
A backtest must use only information that was actually available at the moment
of the simulated decision. Four leaks recur:
\begin{itemize}
  \item \textbf{Restatements.} Fundamentals are revised after publication.
        Using today's corrected value in a simulated decision made years ago
        gives the model knowledge nobody had.
  \item \textbf{Survivorship.} An index constituent list downloaded today
        contains only firms that still exist. Backtesting on it quietly deletes
        every bankruptcy.
  \item \textbf{Look-ahead in features.} A ratio computed from a quarter's
        close, applied to decisions within that quarter.
  \item \textbf{Random splits.} Shuffling time series lets the model train on
        the future (\cref{ch:machine-learning}). Split chronologically, always.
\end{itemize}
Each of these makes the backtest better and the deployment worse. A strategy
that looks excellent and cannot be explained mechanically is usually a leak,
not an edge.
\end{cautionbox}

\begin{cautionbox}[Proxies reintroduce what the law removed]
Removing a protected attribute from the feature list does not remove it from
the model. Postcode encodes ethnicity in many cities; device type and shopping
category encode income; name encodes gender and origin; even application
timestamp can encode shift work. A model trained without the attribute can
reconstruct it and discriminate just as effectively --- while producing a
compliance document stating the attribute was not used.

The test is empirical, not documentary: measure outcome rates across protected
groups on held-out data, and check whether the protected attribute is
predictable from the features you did use. If it is, you have a proxy.
\end{cautionbox}

\begin{practicallab}[Build the backtest you are allowed to believe]
\textbf{Goal.} Produce a model evaluation that would survive a regulator or a
risk committee.

\textbf{Protocol.}
\begin{enumerate}
  \item Take any public time-series dataset with a label you could have known
        at decision time. Write down the exact decision moment.
  \item Split \emph{chronologically} into train, validation, and a final
        holdout period you will touch once.
  \item Establish three baselines: always-negative, a single-feature rule, and
        the previous period's value.
  \item Fit a model. Report precision, recall, and the review-capacity
        implication at three thresholds.
  \item \textbf{Leak hunt.} For every feature, state the timestamp at which its
        value became knowable. Remove anything knowable only later. Re-run and
        record how much performance drops --- that drop was your leak.
  \item Apply realistic costs: transaction fees, review time, or the cost of a
        wrongly declined customer. Recompute whether the model still wins.
  \item Report subgroup outcome rates.
\end{enumerate}

\textbf{Deliverable.} A table with all three baselines, both threshold
analyses, the pre- and post-leak-hunt numbers, and the subgroup rates.
\textbf{The point:} step 5 usually removes most of the apparent performance,
and step 6 removes most of the rest.
\end{practicallab}

\begin{checkpoint}
Your credit model has AUC $0.82$ and the incumbent rule-based system has
$0.79$. Name three things you would need to establish before recommending the
switch, at least one of which is not a performance measurement.
\end{checkpoint}

\chapterreview{Onboard}{Assess}{Transact}{Govern}

\section{Further Reading}
% ============================================================

\subsection*{Further Reading}


\begin{v3topic}{Further Reading}
    \begin{itemize}
        \item \fullcite{jansenMachineLearningAlgorithmic2020} --- The definitive practical reference for ML applied to trading and investment: alpha factor research, portfolio optimization, backtesting, NLP for finance, and deep RL trading agents.
        \item \fullcite{kardellBuildFinancialSoftware2025} --- Hands-on guide to building core banking software with generative AI: ACH processing, FastAPI, PostgreSQL, and AI-assisted development in a FinTech context.
    \end{itemize}

\end{v3topic}
